\section{Abstract SPG Definition and Constraint Domain}
\label{sec:semantic-prefix-grammars}

\subsection{Completability}

Let $L \subseteq \Sigma^*$ be a formal language over a character alphabet $\Sigma$.
A prefix $s \in \Sigma^*$ is \textbf{completable} for $L$ when
\[
	\exists s' \in \Sigma^*, \; ss' \in L.
\]

\subsection{Grammar}
\label{sec:gram-def}

A Semantic Prefix Grammar is a tuple
\[
G = (N, T, P, S, \Theta, \mathcal{T}, \mathcal B, A)
\]
where:
\begin{itemize}[nosep]
    \item \(N\) is a finite set of non-terminals
    \item \(T\) is a finite set of terminal recognizers over segment text
    \item \(P\) is a finite set of productions
    \item \(S \in N\) is the distinguished start symbol
    \item \(\Theta\) is a finite set of semantic rules
    \item \(\mathcal{T} : N \rightharpoonup \Theta\) maps semantic non-terminals to their rule
    \item $\mathcal B$ is the set of binding identifiers
    \item \(A : N \to P^*\) maps each non-terminal \(n\) to its ordered list of productions
    \[
    A(n) = (p_1^n, \dots, p_{k_n}^n).
    \]
\end{itemize}

\begin{assumption}
\label{ass:spg-productivity}
    An \spg has no structural productivity constraints. We will consider the unproductive case trivial in terms of syntactic and semantic constraints, and from now on assume we are dealing with only productive grammars in the \spg class. Thus every reachable alternative in
    the syntactic projection has some terminal yield.
\end{assumption}

\subsubsection{Productions and Alternatives}

\begin{definition}[Production]
A \emph{production} \(p^n_a \in A(n)\) is a finite sequence
\[
    p^n_a
    =
    \alpha^a_1[b^a_1]\,
    \alpha^a_2[b^a_2]\cdots
    \alpha^a_m[b^a_m],
\]
where:
\begin{itemize}[nosep]
    \item \(a \in \{1,\ldots,k_n\}\) is the index of the alternative;
    \item each \(\alpha^a_i \in T \cup N\) is either a terminal recognizer or a non-terminal;
    \item each \(b^a_i \in \mathcal{B} \cup \{\varepsilon\}\) is either a binding identifier or the empty binding.
\end{itemize}
The production \(p^n_a\) is attached to the unique non-terminal \(n\). We write $|p^n_a| = m$ for its length.
\end{definition}

Bindings are syntactic annotations on child positions. They don't affect the context-free structure but they name child evidence that may be referenced by semantic rules.

\paragraph{Segmentation.}
Parsing is parameterized by a segmentation policy
\[
    \mathsf{seg}: \Sigma^* \to \mathsf{Seg}^* .
\]
We require \(\mathsf{seg}\) to be \emph{deterministic and prefix-aware}: an
unfinished final segment may be marked open, but the segmentation of already
closed interior segments is stable under extension. The segmentation algorithm
is a greedy left-to-right longest-match policy. Delimiters are chosen depending
on the task, but independently of the grammar. 

Terminal matching is abstracted through a segment-level interface
\[
    \mathsf{match}(t,x)
    \to
    \{\bot\}
    \cup
    \{\mathsf{Prefix},\mathsf{Exact}, \mathsf{Extensible}\} \times \mathsf{Reg}
\]
This interface is justified by RegEx derivatives ~\cite{brzozowski1964derivatives,owens2009regex} for regular terminals. The parser itself only depends on the evidence exposed by \(\mathsf{match}\). The value \(\bot\) means that the segment cannot match the terminal. 
\[
    \mathsf{Term}(x,\xi, r),
    \qquad
    \xi\in\{\mathsf{Prefix},\mathsf{Exact}, \mathsf{Extensible}\}\qquad r = \partial_x t
\]
Here \(\mathsf{Reg}\) is the set of regular residual recognizers. Terminal evidence is treated uniformly with child evidence: it may satisfy bindings, may remain open at the right frontier, and may be refined by later input.
Here \(r\) is the derivative of terminal recognizer \(t\) after
consuming \(x\). If the derivative is nonempty, the current terminal evidence
still has a completion.


\subsubsection{Binding Signatures}
\label{binding}

Semantic rules refer to children of a production through binding identifiers.
We restrict bindings so that every semantic non-terminal has a fixed binding
interface shared by all of its alternatives.

\begin{definition}[Binding signature]

The \emph{binding signature} for $n \in N$ is the set $\mathcal{B}(n) \subseteq \mathcal B$. 
A grammar is \emph{binding-regular} if for every \(b \in \mathcal B(n)\), $b$ appears only in productions of $n$ and exactly once in every production. 

\[
\begin{aligned}
\forall n \in N,\ \forall b \in \mathcal B(n),\quad
&\Bigl(
  \forall a \in \{1,\cdots,k_n\},
  \exists! i \in \{1,\cdots,|p_a^n|\}.\ b^a_i = b
\Bigr)
\\
&{}\land
\Bigl(
  \forall n' \in N \setminus \{n\},
  \forall a' \in \{1,\cdots,k_{n'}\},~ \forall j \in \{1,\cdots,|p_{a'}^{n'}|\}, \ b^{a'}_j \neq b
\Bigr)
\end{aligned}
\]
\end{definition}

\begin{definition}[Binding position]
Let \(G\) be binding-regular. For \(n \in N\), \(p^n_a \in A(n)\), and
\(b \in \mathcal{B}(n)\), define
\[
    \mathsf{pos}_{n,a}(b)
    =
    \text{the unique } i \in \{1,\ldots,|p^n_a|\}
    \text{ such that } b^a_i = b .
\]
This function is well-defined by totality and uniqueness of binding signatures.
\end{definition}

The parser uses \(\mathsf{pos}_{n,a}\) to route semantic obligations to the
corresponding child. If an item is currently positioned before child \(i\) of
alternative \(p^n_a\), then the only binding obligations projected to that child
are those for bindings \(b\) satisfying
\[
    \mathsf{pos}_{n,a}(b) = i .
\]

\begin{property}
\label{prop:binding-uniqueness}
By definition of $\mathsf{pos}$, and using \emph{binding regularity}, we can establish binding positions are unique :

$$
b_1 \neq b_2 \iff \mathsf{pos}_{n,a}(b_1) \neq \mathsf{pos}_{n,a}(b_2)
$$
\end{property}

\noindent When referring to an SPG in the case of a proof or a definition, we assume binding regularity.

\subsubsection{Syntactic state}
We define a syntactic state that holds our constraint structure, in 
order to prove its correctness.

\begin{definition}[Syntactic State]
For an input $s \in \Sigma^*$, the \emph{syntactic state} is a pair
\[
\sigma(s) = (s, t)
\]
where $t$ is an AST over $G$. We inductively define $t$ as either
\begin{itemize}
  \item a \emph{terminal leaf} $\mathsf{Term}(x, \xi, r)$ with $x \in \mathsf{Seg}$, $\xi \in \{\mathsf{Prefix}, \mathsf{Exact}, \mathsf{Extensible}\}$, and $r \in \Sigma^*$ a witnessing extension;
  \item an \emph{internal node} $\langle n, p, \chi \rangle$ where $n \in N$, $p = \alpha_1^a[b_1^a] \cdots \alpha_m^a[b_m^a] \in A(n)$, and $\chi = (c_1, \ldots, c_k)$ with $k \le m$ is a list of children such that for each $l \in \{1, \ldots, k\}$, $c_l$ is an AST whose root recognizes $\alpha_l^a$.
\end{itemize}
The yield of $t$ — the concatenation of segment text of terminal leaves in left-to-right order — must be consistent with $\mathsf{seg}(s)$.
\end{definition}



The status of an AST node is propagated bottom-up:
\begin{itemize}
  \item a terminal leaf $\mathsf{Term}(x, \xi, r)$ has status $\xi$;
  \item an internal node $\langle n, p, \chi \rangle$ with 
    $\chi = (c_1, \ldots, c_k)$ and $|p| = m$ has status
    \begin{itemize}
      \item $\mathsf{Exact}$ if $k = m$ and $c_k$ is $\mathsf{Exact}$,
      \item $\mathsf{Extensible}$ if $k = m$ and $c_k$ is 
        $\mathsf{Extensible}$,
      \item $\mathsf{Prefix}$ otherwise.
    \end{itemize}
\end{itemize}
The status of $\sigma(s)$ is the status of the root of $t$. An
$\mathsf{Exact}$ state has been fully consumed; an $\mathsf{Extensible}$
state is fully consumed but its rightmost terminal could still match
further input; a $\mathsf{Prefix}$ state has its open frontier on the
rightmost descendant chain of $t$.

\subsubsection{Evidence summary and binding resolution}
\label{sec:evidence-summary}

So far the syntactic state exposes only the AST $t$. We attach to every internal AST node a single \emph{evidence summary}, the canonical node form used throughout the paper: the abstract constraint graph (\S\ref{sec:semantic-domains}) relates such summaries, and the parser arena (\S\ref{sec:engine-parsing}) materializes them.

\begin{definition}[Evidence summary]
\label{def:evidence-summary}
The summary of an internal AST node $\langle n,\,p^n_a,\,\chi\rangle$ at segment span $[i,j)$ is the tuple
\[
  \nu \;=\; \langle n,\;[i,j),\;\rho,\tau, \nabla, \;\beta \rangle ,
\]
where $\rho \in \{\mathsf{Exact},\mathsf{Prefix},\mathsf{Extensible}\}$ is the propagated status, and $\beta$ is the binding-resolution map of Definition~\ref{def:beta-resolution} below. It also carries resolved constrained information in $\tau$ and exported modifications (i.e context operations) in $\nabla$ instantiated by the constraint domain (\S\ref{sec:semantic-domains},~\S\ref{sec:typing-domain}). A terminal leaf has summary $\mathsf{Term}(x,\xi,r)$ as in \S\ref{sec:gram-def} and is treated as a degenerate $\nu$ with $\rho = \xi$.
\end{definition}

\begin{definition}[Frontier-aware binding resolution]
\label{def:beta-resolution}
Let $\nu$ summarize $\langle n,\,p^n_a,\,(c_1,\ldots,c_k)\rangle$, and fix $b \in \mathcal{B}(n)$ with $i_b = \mathsf{pos}_{n,a}(b)$ (\S\ref{binding}). The \emph{binding resolution} of $b$ at $\nu$ is
\[
  \beta(\nu, b) \;=\;
  \begin{cases}
    \mathsf{Resolved}(\nu_{c_{i_b}}) & i_b \le k,\\[0.25em]
    \mathsf{Pending}(i_b)            & i_b > k.
  \end{cases}
\]
$\mathsf{Pending}(i_b)$ records that the right frontier $k$ has not yet reached the binding position of $b$, so no child evidence is available. As $k$ advances on the same item, $\mathsf{Pending}(i_b)$ may transition to $\mathsf{Resolved}$, never the reverse.
\end{definition}


The split $\mathsf{Resolved}/\mathsf{Pending}$ is a concrete instance of the \emph{available}/\emph{unavailable} attribute distinction of incremental attribute grammars defined by Demers et al. ~\cite{demers1981incremental}: an unavailable instance in their works corresponds here to a binding whose position lies past the dot, and its eventual \emph{stabilization} is what the parser obligation store achieves online. The convention also matches the typed-hole formalism of Hazel~\cite{omar2017hazelnut,zhao2024marking}, in which unfilled positions are first-class members of the static state. The $\mathsf{Pending}$ verdict is our parser-time analog of a Hazel hole at a binding-reference vertex.

\subsection{Semantic graphs}
\label{sec:semantic-domains}

We define a semantic layer that records information as a finite \emph{relational graph} over the evidence summaries of \S\ref{sec:evidence-summary}. Constraints operate on a partial (prefix) tree, but the semantic relations themselves do not need to be tree-shaped. Rules relate evidence summaries directly, for example by equality, lookup, unification, subtyping, or more complex decidable predicates. This is the prefix-conservative reading of the \emph{compound dependency graph} of attribute grammars~\cite{knuth1968semantics,demers1981incremental}: vertices are attribute (evidence) instances and edges are semantic dependencies.

\begin{definition}[Evidence graph]
\label{def:evidence-graph}
Let $s \in \Sigma^\ast$ be an input prefix such that $\sigma(s)$ is reachable. The \emph{evidence graph} at $s$ is a finite hypergraph $\mathcal G(s) = (\mathcal E(s),\, \mathcal C(s))$ where:
\begin{itemize}
  \item $\mathcal E(s)$ is the finite set of evidence vertices at $s$. Each vertex is either an evidence summary $\nu = \langle n,[i,j),\rho,\tau,\nabla,\beta\rangle$ (Definition~\ref{def:evidence-summary}), a terminal leaf $\mathsf{Term}(x,\xi,r)$, or a binding-reference $\beta(\nu,b)$ in the sense of Definition~\ref{def:beta-resolution};
  \item $\mathcal C(s) \subseteq \mathcal E(s) \times \mathcal E(s)$ is the finite set of constraint edges emitted so far, drawn from a decidable constraint language and interpreted over $\mathcal E(s)$.
\end{itemize}
\end{definition}

\paragraph{Open and closed graphs.}
The open / closed split is \emph{derived} from frontier evidence rather than asserted in prose. Following the available / unavailable attribute distinction of \cite{demers1981incremental}, a vertex is \emph{frontier} exactly when it still carries some incomplete-input signal:
\[
  \begin{aligned}
  \mathsf{frontier}(\nu) \;&\Longleftrightarrow\; \rho(\nu)\in\{\mathsf{Prefix},\mathsf{Extensible}\}, \\
  \mathsf{frontier}(\beta(\nu,b)) \;&\Longleftrightarrow\; \beta(\nu,b)=\mathsf{Pending}(\cdot), \\
  \mathsf{frontier}(\mathsf{Term}(x,\xi,r)) \;&\Longleftrightarrow\; \xi\in\{\mathsf{Prefix},\mathsf{Extensible}\}.
  \end{aligned}
\]
By construction, we can assert that only one child can be at the frontier for a given prefix state node.
\begin{definition}[Open / closed]
\label{def:open-closed}
$\mathcal G(s)$ is \emph{open} iff some constraint edge $c \in \mathcal C(s)$ is incident to a frontier vertex. $\mathcal G(s)$ is \emph{closed} iff it is not open and contains no contradiction in the underlying constraint domain.
\end{definition}

This semantic notion of closure is distinct from language-level acceptance: a closed evidence graph may still sit inside an incomplete syntactic projection, and (under productivity, Assumption~\ref{ass:spg-productivity}) every open graph admits an extension that closes it.

\begin{example}
\label{ex:open-closed-lambda}
In a typed lambda grammar, the prefix $\lambda x : \mathsf{Int}$ already exposes evidence for the binder $x$ and for its type annotation, but the body binding is still $\mathsf{Pending}$ in the sense of Definition~\ref{def:beta-resolution}: the dot of the lambda item has not yet reached the body position. By Definition~\ref{def:open-closed} the graph is therefore \emph{open}, the constraint edge asserting the body type being incident to a $\mathsf{Pending}$ vertex. After reading a complete annotation and body such as $\lambda x : \mathsf{Int}.\,x$, every binding-reference becomes $\mathsf{Resolved}$ and the corresponding evidence graph is closed.
\end{example}

\begin{definition}[Witness]
\label{def:witness}
Let \(\mathcal G(s) = (\mathcal E(s), \mathcal C(s))\) be the evidence graph at a reachable input prefix \(s\). A \emph{witness} for \(\mathcal G(s)\) is a string \(r \in \Sigma^\ast\) such that
\[
  \mathcal G(s \cdot r) \text{ is closed}.
\]
That is, extending the input by \(r\) drives the evidence graph into a closed satisfying state. The empty string \(\varepsilon\) is a witness exactly when \(\mathcal G(s)\) is already closed.
\end{definition}

\begin{definition}[Denotation]
\label{def:denotation}
The \emph{denotation} of \(\mathcal G(s)\) is the set of its witnesses,
\[
  \mathcal D(\mathcal G(s))
  \;=\;
  \bigl\{\,
    r \in \Sigma^\ast
    \;\big|\;
    \mathcal G(s \cdot r) \text{ is closed }
  \,\bigr\}.
\]
In every realizable instance, there exists a constraint-preserving morphism from
\(\mathcal G(s)\) into a closed semantic graph corresponds to a concrete string
witness, and conversely each witness \(r\) induces such a morphism through the
parser state reached at \(s\cdot r\).
\end{definition}

\begin{definition}[Constraint domain]
\label{def:constraint-domain}
A \emph{constraint domain} is a triple
\[
  D = ( \mathsf{Rules},\, \mathsf{Closed},\, \mathit{eval})
\]
where:
\begin{itemize}
  \item \( \mathsf{Rules}\) is a decidable language of relational rules over evidence nodes and domain constants;
  \item \(\mathsf{Closed}\) is the class of closed evidence graphs admitted by the domain;
  \item \(\mathit{eval}\) is the ideal evaluator,
  \[
    \mathit{eval}(\mathcal G(s))
    \in
    \{\mathsf{Satisfied},\,\mathsf{Live},\,\mathsf{Lost}\}.
  \]
\end{itemize}
It is defined by the following cases:
\[
  \mathit{eval}(\mathcal G(s)) = \mathsf{Satisfied}
  \quad\Longleftrightarrow\quad
  \mathcal G(s) \in \mathsf{Closed}
\]
\[
  \mathit{eval}(\mathcal G(s)) = \mathsf{Live}
  \quad\Longleftrightarrow\quad
  \mathcal G(s) \notin \mathsf{Closed}
  \ \land\
  \mathcal D(\mathcal G(s)) \neq \emptyset
\]
The evaluator returns
\[
  \mathit{eval}(\mathcal G(s)) = \mathsf{Lost}
\]
otherwise, meaning the denotation is empty.
\end{definition}

\begin{lemma}[Safe pruning]
\label{lem:safe-pruning}
If the parser discards an input prefix \(s\) only when
\[
  \mathit{eval}(\mathcal G(s)) = \mathsf{Lost},
\]
then it never discards a prefix that has a closed satisfying continuation.
\end{lemma}

\begin{proof}
Suppose the parser discards \(s\). By assumption,
\[
  \mathit{eval}(\mathcal G(s)) = \mathsf{Lost}.
\]
By Definition~\ref{def:constraint-domain} we have
\(\mathcal D(\mathcal G(s)) = \emptyset\): the prefix \(s\) admits no witness.
Hence discarding \(s\) cannot remove a prefix that could still be completed into a closed satisfying graph.
\end{proof}

\noindent We write \(\mathsf{SPG}(D)\) for grammars over a constraint domain \(D\) equipped with the evaluator of Definition~\ref{def:constraint-domain}. Grammars in this class are said to be realizable, as their semantic constraints match their syntactic realization possibilities.
